\documentclass[10pt,a4paper]{amsart}
\usepackage[margin=19mm]{geometry}
\usepackage{amsmath,amssymb,mathtools}
\usepackage[scr=rsfs,cal=euler]{mathalfa}
\usepackage{xfrac}
\usepackage[unicode,hidelinks,bookmarks=false]{hyperref}
\newcommand{\ie}{i.e.\ }
\newcommand{\cf}{cf.\ }
\newcommand{\resp}{respectively\ }
\newcommand{\Subsection}{\subsection}
\providecommand{\nobreakdash}{\nobreak\hbox{-}}
\setlength{\emergencystretch}{2em}
\multlinegap0pt
\usepackage{bm}
\usepackage{bbm}
\usepackage{dsfont}
\usepackage{xspace}
\usepackage{cancel}
\usepackage{mathtools}
\usepackage{amsthm}
\makeatletter
\@ifundefined{proposition}{\newtheorem{proposition}{Proposition}}{}
\@ifundefined{remark}{\newtheorem{remark}{Remark}}{}
\makeatother
\newcommand{\av}{\mathrm{av}}
\newcommand{\mf}{\mathfrak}
\newcommand{\osc}{\mathrm{osc}}
\newcommand{\pa}{\partial}
\newcommand{\ve}{\varepsilon}
\title[Future stability of solutions of the Einstein-nonlinear scal]{Future stability of solutions of the Einstein-nonlinear scalar field system with decelerated expansion}
\author{Louie Bernhardt}
\date{Source: arXiv:2508.15303v1 (CC BY 4.0)}
\begin{document}
\maketitle

\noindent\emph{Source and licence.} Future stability of solutions of the Einstein-nonlinear scalar field system with decelerated expansion. Version arXiv:2508.15303v1 (CC BY 4.0), distributed under the Creative Commons Attribution 4.0 International licence, as recorded in the arXiv licence metadata for this exact version and shown on the abstract page https://arxiv.org/abs/2508.15303v1. Full rights notice: licenses/arxiv-2508.15303-CC-BY-4.0.txt.\par\smallskip

\noindent\textit{Editorial scope.} Counted unit: the Remark whose source label is \texttt{sec:avgsys.rmk:bottleneck} in the article (source0.tex lines 1314--1316, source label \texttt{sec:avgsys.rmk:bottleneck}), delivered together with the article's own complete original proof of it (source0.tex lines 1317--1420, one \texttt{proof} environment of 104 lines). No mathematics is written, repaired or re-derived: statement and proof are the article's own text line for line. Required context retained uncounted: sec:gaugedeqs.eq:bounddata (lines 916--918); sec:globex.eq:bootstrap (lines 969--971); sec:avgsys.eq:avgbound (lines 1081--1083); sec:avgsys.eq:odeavg2 (lines 1289--1296); sec:avgsys.eq:odeavghori2 (lines 1290--1295); sec:avgsys.eq:odeavglapse2 (lines 1290--1295); sec:avgsys.eq:odeavgscalar2 (lines 1290--1295); sec:avgsys.eq:odeavgshift2 (lines 1290--1295); sec:avgsys.prop:avgest (lines 1308--1313); sec:avgsys.eq:avgest (lines 1310--1312). Every definition, display and label of those lines is retained unchanged. Omitted from this package and not redistributed: 1--915, 919--968, 972--1080, 1084--1288, 1296--1307, 1313--1313, 1421--2160. Nothing inside a retained excerpt is omitted or altered: every retained body is delivered exactly as the article's lines are written, so no omission receipt of any kind accompanies this package. Citations: the retained excerpts carry no citation command at all, so no bibliography entry of the article is transcribed and no reference is redistributed. Reference adaptation: none was needed and none was made; every reference inside the delivered unit resolves inside it, so no unresolved reference or citation is produced. Printed slips kept literal: no printed word, symbol, hypothesis or number of the counted statement or of its proof was altered. Layout and class: the article's own class is \texttt{article} with packages including graphicx, tikz, amsmath, amsfonts, amsthm, amssymb, paralist, xcolor, slashed, accents, esint, mathrsfs, inputenc, titlesec; the wrapper is the lane's pinned \texttt{amsart} editorial wrapper at 10pt on A4, which restates the theorem-like environments the retained text needs and adds the article's own macro definitions that the retained text needs (\texttt{\textbackslash av}, \texttt{\textbackslash mf}, \texttt{\textbackslash osc}, \texttt{\textbackslash pa}, \texttt{\textbackslash ve}), with extra wrapper packages (bm, bbm, dsfont, xspace, cancel, mathtools). The delivered numbering is the unit's own. Assets: no retained span contains a figure environment, a graphics-inclusion command, a PDF-page inclusion, a table or a TikZ picture; no image file is copied into this package, the package declares no asset dependency and no image file is redistributed with it. Licence: version arXiv:2508.15303v1 of this article is distributed under the Creative Commons Attribution 4.0 International licence, as recorded in the arXiv licence metadata for this exact version and shown on the abstract page https://arxiv.org/abs/2508.15303v1. A full rights notice is delivered at licenses/arxiv-2508.15303-CC-BY-4.0.txt. Characters: every retained excerpt is pure ASCII, so the delivered engine, XeLaTeX with its default Latin Modern text font, reports no missing character.
        \begin{equation}\label{sec:gaugedeqs.eq:bounddata}
        S_K(\tau_0) \leq \ve^2,
        \end{equation}

\begin{equation}\label{sec:globex.eq:bootstrap}
    S_K(\tau) \leq \ve
\end{equation}

\begin{equation}\label{sec:avgsys.eq:avgbound}
    |u_\av| \leq C\|u\|_{L^2},\quad \|u_{\osc}\|_{L^2}\leq C\|u\|_{L^2},
\end{equation}

\begin{subequations}\label{sec:avgsys.eq:odeavg2}
    \begin{gather}
        \Big|\pa_\tau  h_{\av}^{00}(\tau) + \frac{4q+2}{\tau}h_{\av}^{00}(\tau)\Big| \leq\> C\Big(\sum_{a,b=1}^3|\pa_\tau h_{\av}^{ab}(\tau)| + \ve^2 \tau^{-2(q-\delta)}\Big),\label{sec:avgsys.eq:odeavglapse2}\\
        \Big|\pa_\tau   h_{\av}^{0i}(\tau) + \frac{2q+2}{\tau}  h_{\av}^{0i}(\tau)\Big| \leq\> C\ve^2\tau^{-2(q-\delta)},\label{sec:avgsys.eq:odeavgshift2}\\
        \Big|\pa_\tau^2 h_{\av}^{ij}(\tau) + \frac{2q+2}{\tau}\pa_\tau  h_{\av}^{ij}(\tau)\Big| \leq C\Big(\frac{1}{\tau^2}|\psi_{\av}(\tau)| + \ve^2\tau^{-2(q-\delta)}\Big),\label{sec:avgsys.eq:odeavghori2}\\
        \Big|\pa_\tau^2 \psi_{\av}(\tau) + \frac{2q+2}{\tau}\pa_\tau \psi_{\av}(\tau) + \frac{2(2q+1)(q+2)}{\tau^2}\psi_{\av}(\tau)\Big| \leq C\ve^2\tau^{-2(q-\delta)}.\label{sec:avgsys.eq:odeavgscalar2}
    \end{gather}
\end{subequations}

\begin{proposition}[Estimate for the spatial averages]\label{sec:avgsys.prop:avgest}
    Suppose the bootstrap assumption \eqref{sec:globex.eq:bootstrap} holds on the interval $I = [\tau_0,T]$, and suppose the initial data satisfies the bound \eqref{sec:gaugedeqs.eq:bounddata}. Then the norm $S_{\av}$ satisfies the following estimate for all $\tau \in I$:
    \begin{equation}\label{sec:avgsys.eq:avgest}
        S_{\av}(\tau) \leq C\ve^2.
    \end{equation}
\end{proposition}

\begin{remark}[Reduced decay of averages from nonlinear coupling]\label{sec:avgsys.rmk:bottleneck}
    The weights in $\tau$ in the various terms that make up the norm $S_{\av}$ rescale the leading-order behaviour of the quantities $(h_{\av}^{\mu\nu},\pa_\tau h_{\av}^{\mu\nu},\psi_{\av},\pa_\tau \psi_{\av})$. Consequently, the improved bound \eqref{sec:avgsys.eq:avgest} on $S_{\av}$ implies the corresponding amount of decay of these quantities. These decay rates are weaker than the principle ODE (the left hand side of the system \eqref{sec:avgsys.eq:odeavg2}) predict. For example, the principle terms in the ODE \eqref{sec:avgsys.eq:odeavgshift2} for the $h_{\av}^{0i}$ imply that $h_{\av}^{0i} = O(\tau^{-(q+1)})$, not $O(\tau^{-q})$ as implied by \eqref{sec:avgsys.eq:avgest}. Ultimately the weaker decay is due to the $\ve^2 \tau^{-2(q-\delta)}$ inhomogeneous terms present on the right hand side of the system \eqref{sec:avgsys.eq:odeavg2}, which arise from the nonlinear coupling of the system to the oscillating terms $(h_{\osc}^{\mu\nu},\psi_{\osc})$. The decay of the oscillating terms limit or bottleneck the decay available from the principle system for the averaged quantities.
\end{remark}

\begin{proof}
    Let us define the following rescaled functions
    \begin{subequations}
        \begin{gather}
            \mf{u}^{00}:= \tau^{q-1}h_{\av}^{00},\qquad \mf{u}^{0i}:= \tau^{q}h_{\av}^{0i},\qquad
            \mf{u}^{ij}:= h_{\av}^{ij},\qquad \mf{v}^{ij}:= \tau^q\pa_\tau\mf{u}^{ij},\\ \mf{r}:= \tau^{q+\delta-1}\psi_{\av},\qquad \mf{s}:=\tau\pa_\tau \mf{r}.
        \end{gather}
    \end{subequations}
    We claim that the system \eqref{sec:avgsys.eq:odeavg2} implies the following differential equalities/inequalities:
    \begin{subequations}
        \begin{gather}
            \frac{1}{2}\pa_\tau (\mf{u}^{00})^2 \leq -\frac{3q+3}{\tau}(\mf{u}^{00})^2 + \frac{C}{\tau}\sum_{a,b=1}^3|\mf{u}^{00}||\mf{v}^{ab}| + C\ve^2\tau^{-(q+1-2\delta)}|\mf{u}^{00}|,\label{sec:avgsys.eq:avgestproof1}\\
            \frac{1}{2}\pa_\tau(\mf{u}^{0i})^2 \leq-\frac{q+2}{\tau}(\mf{u}^{0i})^2 + C\ve^2\tau^{-(q-2\delta)}|\mf{u}^{0i}|,\label{sec:avgsys.eq:avgestproof2}\\
            \frac{1}{2}\pa_\tau(\mf{u}^{ij})^2 = \frac{1}{\tau^{q}}\mf{u}^{ij}\mf{v}^{ij},\qquad \frac{1}{2}\pa_\tau \mf{r}^2 =\frac{1}{\tau}\mf{r}\mf{s}\label{sec:avgsys.eq:avgestproof3}\\
            \frac{1}{2}\pa_\tau (\mf{v}^{ij})^2 \leq - \frac{q+2}{\tau}(\mf{v}^{ij})^2 + \frac{C}{\tau^{1+\delta}}|\mf{v}^{ij}||\mf{r}| + C\ve^2 \tau^{-(q-2\delta)}|\mf{v}^{ij}|,\label{sec:avgsys.eq:avgestproof4}\\
            \frac{1}{2}\pa_\tau \mf{s}^2 \leq -\frac{3-2\delta}{\tau}\mf{s}^2 - \frac{c_1}{\tau}\mf{s}\mf{r} + C\ve^2\tau^{-(q-3\delta)}|\mf{s}|,\label{sec:avgsys.eq:avgestproof5}
        \end{gather}
    \end{subequations}
    where $c_1$ is the positive constant
    \begin{equation}
        c_1 = 3q^2+9q+6-3\delta+\delta^2.
    \end{equation}
    The identities in \eqref{sec:avgsys.eq:avgestproof3}, are true by the definitions of the $\mf{v}^{ij}$ and $\mf{s}$ respectively. The inequality \eqref{sec:avgsys.eq:avgestproof1} follows from the computation
    \begin{align}
        \frac{1}{2}\pa_\tau (\mf{u}^{00})^2 &= \tau^{q-1}\mf{u}^{00}\big(\pa_\tau h_{\av}^{00} + \frac{q-1}{\tau}h_{\av}^{00}\big)\nonumber\\
        &= - \frac{3q+3}{\tau}(\mf{u}^{00})^2 + \tau^{q-1}\mf{u}^{00}\big(\pa_\tau h_{\av}^{00} + \frac{4q+2}{\tau}h_{\av}^{00}\big)\nonumber\\
        &\leq -\frac{3q+3}{\tau}(\mf{u}^{00})^2 + \frac{C}{\tau}\sum_{a,b=1}^3|\mf{u}^{00}||\pa_\tau h_{\av}^{ab}| + C\ve^2\tau^{-(q+1-2\delta)}|\mf{u}^{00}|,
    \end{align}
    where in the last line we used \eqref{sec:avgsys.eq:odeavglapse2}. The inequality \eqref{sec:avgsys.eq:avgestproof2} follows similarly from \eqref{sec:avgsys.eq:odeavgshift2}. For \eqref{sec:avgsys.eq:avgestproof4}, we have
    \begin{align}
        \frac{1}{2}\pa_\tau (\mf{v}^{ij})^2 =&\>\tau^q\mf{v}^{ij}\big(\pa_\tau^2 h_{\av}^{ij} + \frac{q}{\tau}\pa_\tau h_{\av}^{ij}\big)\nonumber\\
        =&\>\tau^q\mf{v}^{ij}\big(\pa_\tau^2h_{\av}^{ij}+\frac{2q+2}{\tau}\pa_\tau h_{\av}^{ij}\big) - \frac{q+2}{\tau}(\mf{v}^{ij})^2\nonumber\\
        \leq&\> \frac{C}{\tau}|\mf{v}^{ij}||\mf{r}| + C\ve^2 \tau^{-(q-2\delta)}|\mf{v}^{ij}|- \frac{q+2}{\tau}(\mf{v}^{ij})^2,
    \end{align}
    where in the last line we used \eqref{sec:avgsys.eq:odeavghori2}. Finally for \eqref{sec:avgsys.eq:avgestproof5}, we compute that
    \begin{align}
        \frac{1}{2}\pa_\tau \mf{s}^2 = &\>\mf{s}\pa_\tau (\tau^{q+\delta} \pa_\tau \psi_{\av} + (q+\delta-1)\tau^{q+\delta-1}\psi_{\av})\nonumber\\
        =&\>\tau^{q+\delta} \mf{s}\Big(\pa_\tau^2\psi_{\av} + \frac{2q+2\delta-1}{\tau}\pa_\tau \psi_{\av} + \frac{(q+\delta-1)^2}{\tau^2}\psi_{\av}\Big)\nonumber\\
        =&\>\tau^{q+\delta}\mf{s}\Big(\pa_\tau^2\psi_{\av} + \frac{2q+2}{\tau}\pa_\tau \psi_{\av} + \frac{2(2q+1)(q+2)}{\tau^2}\psi_{\av}\Big)\nonumber\\
        &-(3-2\delta)\tau^{q+\delta-1}\mf{s}\pa_\tau \psi_{\av} - \big(3q^2+(12-2\delta)q+(3+2\delta-\delta^2)\big)\tau^{q+\delta-2}\mf{s}\psi_{\av}.
    \end{align}
    We estimate by \eqref{sec:avgsys.eq:odeavgscalar2}:
    \begin{equation}
        \tau^{q+\delta}\mf{s}\Big(\pa_\tau^2\psi_{\av} + \frac{2q+2}{\tau}\pa_\tau \psi_{\av} + \frac{2(2q+1)(q+2)}{\tau^2}\psi_{\av}\Big) \leq C\ve^2\tau^{-(q-3\delta)}|\mf{s}|,
    \end{equation}
    moreover
    \begin{multline}
        -(3-2\delta)\tau^{q+\delta-1}\mf{s}\pa_\tau \psi_{\av} - \Big[3q^2+(12-2\delta)q+(3+2\delta-\delta^2)\Big]\tau^{q+\delta-2}\mf{s}\psi_{\av}\\
        = - \frac{3-2\delta}{\tau}\mf{s}^2 - \frac{3q^2+9q+6-3\delta+\delta^2}{\tau^2}\mf{s}\mf{r}.
    \end{multline}
    The constant $c_1$ is clearly positive for $\delta$ sufficiently small. Combining the above computations, \eqref{sec:avgsys.eq:avgestproof5} follows.
    
    Next we derive estimates for the system $(\mf{u}^{\mu\nu}, \mf{v}^{ij}, \mf{r},\mf{s})$ using the system \eqref{sec:avgsys.eq:avgestproof1} - \eqref{sec:avgsys.eq:avgestproof5}. The estimate
    \begin{equation}\label{sec:avgsys.eq:avgestproof7}
     (\mf{u}^{0i}(\tau))^2 \leq C\big((\mf{u}^{0i}(\tau_0))^2 + \ve^4\tau^{-\delta}\big)
    \end{equation}
    follows from \eqref{sec:avgsys.eq:avgestproof2} by a straightforward application of Cauchy-Schwartz and Gr\"onwall's inequality. For the remaining quantities, we estimate the $\mf{r}$, $\mf{s}$ (which are derived from $\psi_{\av}$) first, since these equations do not contain mixed linear terms, while the equations for the $h_{\av}^{00}$, $h_{\av}^{ij}$ quantities contain linear terms that depend on $\psi_{\av}$. We compute from \eqref{sec:avgsys.eq:avgestproof3} and \eqref{sec:avgsys.eq:avgestproof5} the inequality
    \begin{equation}
        \pa_\tau\big(\frac{c_1}{2}\mf{r}^2 + \frac{1}{2}\mf{s}^2\big) \leq \frac{c_1}{\tau}\mf{r}\mf{s} -\frac{3-2\delta}{\tau}\mf{s}^2 - \frac{c_1}{\tau}\mf{r}\mf{s} + C\ve^2 \tau^{-q + 3\delta}|\mf{s}| \leq C\ve^2\tau^{-q+3\delta}|\mf{s}|.
    \end{equation}
    Note the cancellation of the mixed $c_1\tau^{-1}\mf{r}\mf{s}$ terms. This implies by a Gr\"onwall-type inequality that
    \begin{equation}\label{sec:avgsys.eq:avgestproof8}
        \frac{c_1}{2}\mf{r}^2(\tau) + \frac{1}{2}\mf{s}^2(\tau) \leq C(\mf{r}^2(\tau_0) + \mf{s}^2(\tau_0) + \ve^4 \tau^{-\delta}),
    \end{equation}
    as long as $\delta$ is chosen sufficiently small that $q \geq 1+4\delta$. We must estimate the $\mf{u}^{00}$ and $\mf{v}^{ij}$ jointly, due to the mixed $|\mf{u}^{00}||\mf{v}^{ab}|$ term present in \eqref{sec:avgsys.eq:avgestproof1}. Letting $c_2$ denote a positive constant to be set in a moment, we compute
    \begin{align}
        \pa_\tau\Big[\frac{1}{2}(\mf{u}^{00})^2 + \frac{c_2}{2} \sum_{a,b=1}^3(\mf{v}^{ab})^2\Big] \leq &\> -\frac{3q+3}{\tau}(\mf{u}^{00})^2 + \frac{C}{\tau}\sum_{a,b=1}^3 |\mf{u}^{00}||\mf{v}^{ab}| - \frac{c_2(q+2)}{\tau}\sum_{a,b=1}^3(\mf{v}^{ab})^2\nonumber\\
        &+ C\ve^2\tau^{-(q+1-2\delta)}|\mf{u}^{00}| + C\sum_{a,b=1}^3\Big(\frac{1}{\tau^{1+\delta}}|\mf{r}||\mf{v}^{ab}| + \ve^2\tau^{-(q-2\delta)}|\mf{v}^{ab}|\Big).
    \end{align}
    We fix $c_2$ to be sufficiently large that the quadratic terms on the RHS of the above equation are non-positive, i.e. we have
    \begin{equation}
        -\frac{3q+3}{\tau}(\mf{u}^{00})^2 + \frac{C}{\tau}\sum_{a,b=1}^3 |\mf{u}^{00}||\mf{v}^{ab}| - \frac{c_2(q+2)}{\tau}\sum_{a,b=1}^3(\mf{v}^{ab})^2 \leq 0.
    \end{equation}
    Then we apply the inequality \eqref{sec:avgsys.eq:avgestproof8} to estimate $|\mf{r}|$ and apply Cauchy-Schwartz to obtain
    \begin{equation}
        \pa_\tau\Big(\frac{1}{2}(\mf{u}^{00})^2 + \frac{c_2}{2} \sum_{a,b=1}^3(\mf{v}^{ab})^2\Big) \leq \frac{C}{\tau^{1+\delta}}\Big((\mf{u}^{00})^2+\sum_{a,b=1}^3(\mf{v}^{ab})^2\Big) + \frac{C}{\tau^{1+\delta}}(\mf{r}^2(\tau_0)+\mf{s}^2(\tau_0)+\ve^4),
    \end{equation}
    from which we obtain
    \begin{equation}\label{sec:avgsys.eq:avgestproof9}
        \frac{1}{2}(\mf{u}^{00}(\tau))^2 + \frac{c_2}{2} \sum_{a,b=1}^3(\mf{v}^{ab}(\tau))^2 \leq C\Big((\mf{u}^{00}(\tau_0))^2 + \sum_{a,b=1}^3(\mf{v}^{ab}(\tau_0))^2 + \mf{r}^2(\tau_0) +\mf{s}^2(\tau_0) + \ve^4\tau^{-\delta}\Big).
    \end{equation}
    by Gr\"onwall. We then derive the bound
    \begin{equation}\label{sec:avgsys.eq:avgestproof10}
        (\mf{u}^{ij})^2 \leq C\Big((\mf{u}^{00}(\tau_0))^2 + \sum_{a,b=1}^3\big[(\mf{u}^{ab})^2+(\mf{v}^{ab}(\tau_0))^2\big] + \mf{r}^2(\tau_0) +\mf{s}^2(\tau_0) + \ve^4\tau^{-\delta}\Big)
    \end{equation}
    from \eqref{sec:avgsys.eq:avgestproof3} through a straightforward application of Cauchy-Schwartz and Gr\"onwall. We convert the above estimates for $(\mf{u}^{\mu\nu},\mf{v}^{\mu\nu},\mf{r},\mf{s})$ into a bound on $S_{\av}(\tau)$ using
    \begin{subequations}
        \begin{gather}
            \tau^{q-1-\delta}| h_{\av}^{00}| \sim \tau^{-\delta}|\mf{u}^{00}|,\\
            \tau^{q-\delta}|h_{\av}^{0i}| \sim |\mf{u}^{0i}|,\\
            |h_{\av}^{ij}| + \tau^q|\pa_\tau h_{\av}^{ij}| \sim |\mf{u}^{ij}| + |\mf{v}^{ij}|,\\
            \tau^{q+\delta-1}|\psi_{\av}| + \tau^{q+\delta}|\pa_\tau \psi_{\av}|  \sim  |\mf{r}| + |\mf{s}|,
        \end{gather}
    \end{subequations}
    and applying the estimates \eqref{sec:avgsys.eq:avgestproof7} - \eqref{sec:avgsys.eq:avgestproof10}. The quantities $\tau^q |\pa_\tau h_{\av}^{00}|$ and $\tau^{q+1}|\pa_\tau h_{\av}^{0i}|$ can be estimated by taking the inequalities \eqref{sec:avgsys.eq:odeavglapse2}, \eqref{sec:avgsys.eq:odeavgshift2} and applying the estimates \eqref{sec:avgsys.eq:avgestproof7} - \eqref{sec:avgsys.eq:avgestproof10} to the right hand side. It follows that
    \begin{equation}
        S_{\av}(\tau) \leq C(S_{\av}(\tau_0) + \ve^2\tau^{-\delta/2}).
    \end{equation}
    The average of a function satisfies the bound \eqref{sec:avgsys.eq:avgbound} and hence the initial data bound \eqref{sec:gaugedeqs.eq:bounddata} implies
    \begin{equation}
        S_{\av}(\tau_0) \leq C\sum_{\alpha,\beta=0}^3\Big(\|h^{\alpha\beta}(\tau_0)\|_{L^2} + \|\pa_\tau h^{\alpha\beta}(\tau_0)\|_{L^2}\Big) + C\big(\|\psi(\tau_0)\|_{L^2} + \|\pa_\tau \psi(\tau_0)\|_{L^2}\big) \leq C\ve^2.
    \end{equation}
    This completes the proof of Proposition~\ref{sec:avgsys.prop:avgest}.
\end{proof}

\end{document}
